% Bagian: sets-functions-relations
% Bab: size-of-sets
% Subbagian: introduction

\documentclass[../../../include/open-logic-section]{subfiles}

\begin{document}

\olfileid[id]{sfr}{siz}{int}

\olsection{Pendahuluan}

Ketika Georg Cantor mengembangkan teori himpunan pada dasawarsa 1870-an,
salah satu tujuannya ialah membuat gagasan tentang suatu kumpulan tak hingga
menjadi lebih dapat diterima---suatu ketakhinggaan aktual, sebagaimana
dikatakan para pemikir Abad Pertengahan. Bagian penting dari upayanya ialah
perlakuannya terhadap \emph{ukuran} berbagai himpunan. Jika $a$, $b$, dan $c$
semuanya berbeda, maka secara intuitif himpunan $\{a, b, c\}$ \emph{lebih
besar} daripada $\{a, b\}$. Namun, bagaimana dengan himpunan tak hingga?
Apakah semuanya sama besar? Ternyata tidak.

Gagasan penting yang pertama di sini ialah enumerasi. Kita dapat
mencantumkan setiap himpunan hingga dengan mencantumkan semua !!{element}snya.
Untuk beberapa himpunan tak hingga, kita juga dapat mencantumkan semua
!!{element}snya jika daftarnya sendiri boleh tak hingga. Himpunan semacam itu
disebut !!{enumerable}. Hasil mengejutkan Cantor, yang akan kita pahami
sepenuhnya pada akhir bab ini, ialah bahwa beberapa himpunan tak hingga tidak
!!{enumerable}.

\end{document}
